\documentclass[a4paper,11pt]{article}

\usepackage{kpfonts}
\usepackage[utf8]{inputenc}
\usepackage[margin=3cm]{geometry}

\usepackage{graphicx}
\usepackage{subcaption}
\usepackage{amsmath, amssymb, amsthm}
\usepackage{mathtools}
\DeclarePairedDelimiter\ceil{\lceil}{\rceil}
\DeclarePairedDelimiter\floor{\lfloor}{\rfloor}
\usepackage{microtype}
\usepackage{booktabs}
\usepackage{enumerate,paralist}
\usepackage{xspace}
\usepackage{comment}
\usepackage{xcolor}
\usepackage{array}
\usepackage[nocompress]{cite}
\usepackage[pdfpagelabels,colorlinks,allcolors=blue]{hyperref}
\usepackage[nameinlink,capitalize,noabbrev]{cleveref}

\theoremstyle{plain}
\newtheorem{definition}{Definition}
\newtheorem{theorem}{Theorem}
\newtheorem{lemma}{Lemma}
\newtheorem{claim}{Claim}
\crefname{claim}{Claim}{Claims}
\Crefname{claim}{Claim}{Claims}
\newtheorem{proposition}{Proposition}
\newtheorem{observation}{Observation}
\theoremstyle{definition}  % switches off use of italics in following env's

\title{Smallest Cubic Non-1-Planar Graph}

%\author{Sergey Pupyrev}{Menlo Park, CA, USA}{spupyrev@gmail.com}{https://orcid.org/0000-0003-4089-673X}{}

%\authorrunning{S. Pupyrev}
%
%\Copyright{Sergey Pupyrev}
%
%\ccsdesc[500]{Theory of computation~Design and analysis of algorithms}
%\ccsdesc[500]{Mathematics of computing~Graph theory}
%\keywords{beyond planarity, 1-planar graph, cubic graph}

%\acknowledgements{}
%
%\hideLIPIcs
%\nolinenumbers

\graphicspath{{pics/}}

\definecolor{defblue}{rgb}{0.121,0.47,0.705}
\definecolor{defred}{rgb}{0.7,0.37,0.37}
\definecolor{defgreen}{rgb}{0.01,0.65,0.20}

\newcommand{\df}[1]{\textcolor{defblue}{\emph{#1}}}
\newcolumntype{C}[1]{>{\centering\arraybackslash}p{#1}}
\newcommand{\yes}{{\it \textcolor{defgreen}{yes}}}
\newcommand{\no}{{\it \textcolor{defred}{no}}}

\begin{document}

\author{
    Sergey Pupyrev \\
    spupyrev@gmail.com
}

\date{}
\maketitle

\begin{abstract}
    A graph
    is 1-planar if it can be drawn in the plane such that each edge is crossed at most
    once---a natural generalization of planar graphs that has received increasing
    attention in graph drawing and beyond-planar graph theory.

    We search for the smallest non-1-planar graph.
\end{abstract}

\section{Introduction}

Arguably among the most natural applications of a tool for recognizing 1-planar graphs is
finding the smallest non-1-planar instance in a family of graphs.
A particularly challenging variant of the problem is to
find the smallest 3-regular (cubic) graph that is not 1-planar. While it is easy to show
(see e.g., \cite{DKMW08}) that large random cubic graphs are not 1-planar, the
smallest such instance is unknown.

Throughout the paper, all graphs are finite, simple, and undirected.  The
\emph{order} of a graph $G$ is $|V(G)|$, and we write $n=|V(G)|$.  A graph is
\emph{cubic} if every vertex has degree three, and it is \emph{subcubic} if
every vertex has degree at most three.  A connected graph is
\emph{biconnected} if it has no cut vertex.  A graph is \emph{planar} if it
has a drawing without crossings.  The \emph{girth} of a graph containing a
cycle is the length of its shortest cycle.

\newtheorem*{problem1b}{Problem 1}
\begin{problem1b}[\cite{11011110_cubic,mathse_cubic,oops}]
    What is the smallest cubic non-1-planar graph?
\end{problem1b}

In an attempt to answer the question, Eppstein~\cite{11011110_cubic} manually constructed 1-planar
drawings of several graphs and suggested that either the Coxeter graph or the Tutte--Coxeter graph
is not 1-planar. As pointed out in \cite{oops}, the former is indeed 1-planar.
The latter is conjectured to be non-1-planar in \cite{oops}.

\section{Lower Bound}

% TODO: reduction rules for 3 faces?

We verify the 1-planarity of several families of cubic graphs; refer
to \cref{table:cubic}. In particular, (i) all cubic graphs with up to $n=24$ vertices and (ii) all cubic bipartite graphs with up to $n=28$ vertices are 1-planar.
Note that there are around $150{\times}10^6$ instances across (i) and (ii),
and the total computation took $\approx$$1.5$ machine-months, averaging less than a second per instance. Thus, verifying the 1-planarity of all cubic graphs with $n<30$ remains a challenging task that will likely require novel ideas.

\begin{table}[!h]
    \small
    \centering
    \caption{Computational results for cubic graphs.  The reported elapsed
    times were obtained by parallel computation.}
    \label{table:cubic}
    \begin{tabular}{|r|r|r|r||C{50pt}|}
        \toprule
        graph class & total count & total runtime & runtime per instance & all $1$-planar? \\
        \midrule
        $n=22$                    & $7,\!319,\!447$ & 5h 20m & 85 ms & \yes \\
        $n=24$                    & $117,\!940,\!535$ & 6d 9h 40m & $150$ ms & \yes \\
        \cmidrule{1-1}
        bipartite $n=26$             & $245,\!627$ & 25m & $200$ ms & \yes \\
        %girth$\ge 6$ $n=26$          & $181,\!227$ & $181,\!227$ & 4130 sec & $700$ ms & \yes \\
        girth$\ge 5$ $n=26$          & $31,\!478,\!584$ & 3d 6h 20m & $290$ ms & \yes \\
        $n=26$                    & $2,\!094,\!480,\!864$ & & & ? \\
        \cmidrule{1-1}
        bipartite $n=28$          & $2,\!291,\!589$ & 6h 55m & $340$ ms & \yes \\
        girth$\ge 6$ $n=28$          & $4,\!624,\!501$ & 2d 23h 20m & $1700$ ms & \yes \\
        $n=28$                    & $40,\!497,\!138,\!011$ & & & ? \\
        \cmidrule{1-1}
        bipartite $n=30$          & $23,\!466,\!857$ & & & ? \\
        \bottomrule
    \end{tabular}
\end{table}

The next table gives corresponding counts for several subcubic graph
families.

\begin{table}[!h]
    \small
    \centering
    \caption{Computational results for selected subcubic graph families.}
    \label{table:subcubic}
    \begin{tabular}{@{}r p{205pt} r C{55pt}@{}}
        \toprule
        graph size & graph family & total count & all $1$-planar? \\
        \midrule
        $n=18$ & connected cubic                         & $41,\!301$ & \yes \\
        $n=18$ & biconnected, nonplanar, triangle-free cubic & $7,\!734$ & ? \\
        \cmidrule{1-1}
        $n=20$ & connected cubic                         & $510,\!489$ & \yes \\
        $n=20$ & biconnected, nonplanar, triangle-free cubic & $97,\!141$ & ? \\
        \cmidrule{1-1}
        $n=22$ & connected cubic                         & $7,\!319,\!447$ & \yes \\
        $n=22$ & connected, triangle-free cubic          & $1,\!435,\!720$ & ? \\
        $n=22$ & biconnected, nonplanar, triangle-free cubic & $1,\!432,\!712$ & ? \\
        \cmidrule{1-1}
        $n=24$ & connected cubic                         & $117,\!940,\!535$ & \yes \\
        $n=24$ & biconnected, nonplanar, triangle-free cubic & $23,\!751,\!969$ & ? \\
        \bottomrule
    \end{tabular}
\end{table}

Our goal is to find the \df{smallest subcubic non-1-planar} graphs.  We order
graphs lexicographically by $(|V|,|E|)$ and let $G$ be a minimum non-1-planar
graph of maximum degree at most three.  Thus every smaller subcubic graph is
1-planar.  This formulation is useful even though the graphs in our
enumeration are cubic: any cubic counterexample contains such a minimal
subcubic counterexample of no larger order.

\begin{lemma}
    A smallest subcubic non-1-planar graph is biconnected.
\end{lemma}

\begin{proof}
    A disconnected graph is 1-planar if each of its components is 1-planar.
    Likewise, 1-planar drawings of the blocks incident to a cut vertex can be
    placed in disjoint wedges around that vertex.  Hence a non-1-planar graph
    has a non-1-planar block.  Such a block would be smaller than $G$ unless
    $G$ itself is biconnected.
\end{proof}

\begin{lemma}
    A smallest subcubic non-1-planar graph is cubic.
\end{lemma}

\begin{proof}
    A vertex of degree zero or one can be deleted and then reinserted without
    creating a crossing.  Assume that $G$ contains a vertex $v$ of degree two,
    with neighbors $u$ and $w$. Consider three cases.
    \begin{itemize}
        \item If $uw\notin E$, remove $v$ and insert $uw$.  The result $H$ is
        a smaller subcubic graph and hence is 1-planar.  In a drawing of $H$,
        replace the curve for $uw$ by the path $(u,v,w)$.  At most one path
        edge inherits the possible crossing of $uw$.

        \item Suppose $uw\in E$ and $u,w$ have another common neighbor $x$.
        If there are vertices outside $\{u,v,w,x\}$, then $x$ is a cut vertex.
        Otherwise the graph is planar.  Both alternatives contradict the
        choice of $G$.

        \item Suppose $uw\in E$ and $u,w$ have no other common neighbor.
        Merge $u,v,w$ into one vertex.  No loop or parallel edge is created,
        and the new vertex has degree at most two.  The result is a smaller
        subcubic graph.  In its 1-planar drawing, expand the merged vertex
        inside a sufficiently small disk into the triangle $(u,v,w)$ while
        preserving all external incidences.
    \end{itemize}
\end{proof}

We use the following elementary replacement observation.

\begin{lemma}[Planar two-terminal replacement]\label{lem:two-terminal-replacement}
    Let $P$ be a planar graph with two vertices $a,b$ incident to a common
    face.  Suppose that $G$ is obtained from a graph $H$ by deleting an edge
    $xy$, inserting a disjoint copy of $P$, and adding the edges $xa$ and
    $yb$.  If $H$ is 1-planar, then $G$ is 1-planar.
    Moreover, let $F$ be a set of edges of $G$.  Form $F_H$ by retaining every
    edge of $F$ outside the replacement and, if $F$ meets the inserted patch
    or either attachment, adding $xy$.  Then $|F_H|\le |F|$, and a 1-planar
    drawing of $H$ in which every edge of $F_H$ is uncrossed expands to one of
    $G$ in which every edge of $F$ is uncrossed.
\end{lemma}

\begin{proof}
    Take a 1-planar drawing of $H$ and remove the curve for $xy$.  Draw $P$
    in a small disk around $x$, using a common face of $a,b$ as its outer
    face.  Draw $xa$ inside the disk and route $yb$ along the former curve of
    $xy$.  All new edges in the disk are crossing-free, while $yb$ inherits
    the at most one crossing of $xy$.
    For the strengthened statement, if $F$ meets the replacement then $xy$ is
    uncrossed.  Choose the routing disk around its whole curve; the planar
    patch and both attachments are then crossing-free.  Every retained
    outside edge keeps its old crossing status.
\end{proof}

\begin{lemma}\label{lem:minimal-triangle-free}
    A smallest cubic non-1-planar graph is triangle-free.
\end{lemma}

\begin{proof}
    For a contradiction, assume $G$ has a triangle $T=(u,v,w)$.  Since $G$ is
    cubic, each vertex of $T$ has exactly one neighbor outside $T$; denote
    them by $u',v',w'$. Consider cases.

    \begin{itemize}
        \item If $u',v',w'$ are pairwise distinct, contract $T$ into one
        vertex.  The result is a smaller simple subcubic graph and is therefore
        1-planar.  Expanding the contracted vertex into $T$ inside a small
        disk introduces no crossing.

        \item If $u'=v'=w'$, then connectedness and cubicity imply that
        $G=K_4$, which is planar.

        \item It remains to consider the case $u'=v'=x$ and $w'=y\ne x$.
        The subgraph $P_0$ induced by $u,v,w,x$ is a diamond.  Its terminals
        are $x,w$, and $P_0+xw=K_4$ is planar.  Let $p$ be the third neighbor
        of $x$, outside the diamond.  If $p=y$, then $y$ is a cut vertex,
        contradicting biconnectivity.  If $p\ne y$ and $py\notin E$, replace
        the diamond by the edge $py$.  This gives a smaller subcubic graph,
        and \cref{lem:two-terminal-replacement} expands any of its 1-planar
        drawings back into a drawing of $G$.

        Finally suppose $py\in E$.  Absorb $p,y$ and the three edges
        $xp,py,yw$ into the planar patch.  The new terminals $p,y$ are
        incident to a common face (indeed, they are adjacent).  Let their two
        as-yet unaccounted-for neighbors be $p_1,y_1$.  If $p_1=y_1$, that
        vertex is a cut vertex.  If they are distinct and nonadjacent, replace
        the enlarged patch by $p_1y_1$ and apply
        \cref{lem:two-terminal-replacement}.  If
        they are adjacent, absorb them as well and repeat.

        Every absorbed terminal pair consists of new vertices: all vertices
        of an earlier pair already have their three incident edges inside the
        patch or leading to the next pair.  Since $G$ is finite, the process
        must therefore reach either a common outside neighbor (a cut vertex)
        or two distinct nonadjacent outside neighbors (a valid replacement).
        Both outcomes contradict the choice of $G$.
    \end{itemize}
\end{proof}

For the flexibility statements we use a strengthened form of the same
induction.
Call a graph \emph{$k$-flexible} if, for every set of at most $k$ of its
edges, it has a 1-planar drawing in which all those edges are uncrossed.

\begin{lemma}[Minimal obstruction to flexibility]\label{lem:flexibility-core}
    Let $k$ be a nonnegative integer.  A lexicographically smallest subcubic graph that is not
    $k$-flexible is biconnected, cubic, and triangle-free.
\end{lemma}

\begin{proof}
    Let $G$ be such a graph.  For any set $F\subseteq E(G)$ with $|F|\leq k$,
    distribute the edges of $F$ among the components or blocks of $G$.  Every
    proper component or block is smaller than $G$ and is therefore
    $k$-flexible.  Its drawing may consequently be chosen with its part of
    $F$ uncrossed.  Placing component drawings in disjoint regions, or block
    drawings in disjoint wedges at their common cut vertices, gives the
    required drawing of $G$.  Hence $G$ is biconnected.

    It remains to exclude a vertex $v$ of degree two.  Let its neighbors be
    $u$ and $w$.  If $uw\notin E(G)$, delete $v$ and add $uw$.  Retain every
    prescribed edge outside $uv,vw$, and prescribe $uw$ if $uv$ or $vw$
    belongs to $F$.  Subdividing an uncrossed drawing of $uw$ restores both
    $uv$ and $vw$ without crossings.

    Suppose instead that $uw\in E(G)$.  If $u$ and $w$ have another common
    neighbor $x$, then either $G$ is planar or $x$ is a cut vertex, contrary
    to the preceding paragraph.  Otherwise identify $u,v,w$ to a single
    vertex.  The resulting graph is simple and subcubic.  Retain prescribed
    edges outside the triangle, map its prescribed attachment edges to their
    images, and omit its prescribed internal edges.  Expanding the identified
    vertex inside a disk restores the triangle, keeps its internal edges
    uncrossed, and preserves the crossing status of every mapped attachment.
    Thus $G$ is cubic.

    Finally, let $T$ be a triangle.  If its three external neighbors are
    distinct, contract $T$ to $z$.  Map a prescribed attachment $v_ix_i$ to
    $zx_i$, retain prescribed outside edges, and omit prescribed triangle
    edges from the mapped set; the map does not increase its size.  Expanding
    an uncrossed mapped attachment and drawing the triangle inside the vertex
    disk preserves every prescribed edge.  In the repeated-neighbor cases,
    apply the iterative two-terminal replacement used in the proof of
    \cref{lem:minimal-triangle-free}.  The prescribed-edge construction in
    \cref{lem:two-terminal-replacement} does not increase the number of
    prescribed edges.  Every case therefore reduces to a smaller subcubic
    graph, a contradiction.
\end{proof}

\begin{lemma}[Drawing criterion for flexibility]\label{lem:drawing-coverage}
    Let $\mathcal D$ be a collection of 1-planar drawings of a graph
    $G$.  For each $D\in\mathcal D$, let $X_D\subseteq E(G)$ contain every
    edge crossed in $D$.  If, for every $F\subseteq E(G)$ with $|F|\leq k$,
    there exists $D\in\mathcal D$ such that $F\cap X_D=\varnothing$, then
    $G$ is $k$-flexible.
\end{lemma}

\begin{proof}
    Fix $F\subseteq E(G)$ with $|F|\leq k$.  By hypothesis, some
    $D\in\mathcal D$ satisfies $F\cap X_D=\varnothing$.  Since $X_D$ contains
    every crossed edge of $D$, every edge of $F$ is uncrossed in $D$.
\end{proof}

To establish $k$-flexibility, we construct $\mathcal D$ iteratively.  If some
set $F\subseteq E(G)$ with $|F|\leq k$ intersects $X_D$ for every
$D\in\mathcal D$, we add a 1-planar drawing in which every edge of $F$ is
uncrossed.  If no such $F$ exists, \cref{lem:drawing-coverage} applies.  The
procedure terminates because the drawing added in an iteration has a
crossed-edge set disjoint from the corresponding set $F$; hence that $F$
cannot occur again.  Only finitely many subsets of $E(G)$ have size at most
$k$.

The following local observation will be used when prescribing uncrossed
edges.

\begin{lemma}[Local crossing elimination]\label{lem:local-crossing-elimination}
    Let $v$ have neighbors $x,y,z$ in a 1-planar drawing, and suppose that
    $vx$ crosses $yz$.  The drawing can be modified so that this crossing is
    removed and no new crossing is introduced.  In particular, every edge
    that was uncrossed before the modification remains uncrossed.
\end{lemma}

\begin{proof}
    Move $v$ into a sufficiently small disk beside the crossing, on the
    $x$-side of $yz$.  Draw $vx$ along the crossing-free portion of its original
    curve, and draw $vy$ and $vz$ in narrow neighborhoods of the two
    crossing-free portions of $yz$.  The modification removes the crossing
    between $vx$ and $yz$ and introduces none.
\end{proof}

\begin{lemma}[4-cycle expansion]\label{lem:4-cycle-expansion}
    Let $G$ be a triangle-free cubic graph containing the cycle
    $C=(v_0,v_1,v_2,v_3)$, and let $x_i$ be the neighbor of $v_i$ outside
    $C$.  Form $H$ by deleting $C$, adding adjacent vertices $a,b$, joining
    $a$ to $x_0,x_1$, and joining $b$ to $x_2,x_3$.  If $H$ has a 1-planar
    drawing in which $ab$ is uncrossed, then $G$ is 1-planar.
\end{lemma}

\begin{proof}
    Triangle-freeness implies $x_0\ne x_1$ and $x_2\ne x_3$, so $H$ is a
    simple cubic graph; equality between opposite external neighbors does not
    create parallel edges.  Take a narrow disk containing the curve $ab$ and small
    neighborhoods of its endpoints, but no crossing.  Inside the disk, split
    $a$ into $v_0,v_1$, split $b$ into $v_2,v_3$, preserve the four external
    incidences, and draw the edges of $C$.  Depending on the rotations at
    $a,b$, the two remaining cycle edges can be drawn either without a
    crossing or crossing each other once.  Every edge inherited from $H$
    retains its crossings, and each new edge has at most one crossing.
\end{proof}

The analogous construction pairing $x_1,x_2$ at $a$ and $x_3,x_0$ at $b$
gives the corresponding second expansion.

\begin{lemma}[Adjacent-edge 4-cycle expansion]\label{lem:adjacent-4-cycle}
    In the setting of \cref{lem:4-cycle-expansion}, let
    $c_i=v_iv_{i+1}$ and $s_i=v_ix_i$, with indices modulo four.  If $H$ has
    a 1-planar drawing in which $ab$ and the image of $s_1$ are uncrossed,
    then $G$ has a 1-planar drawing in which the three consecutive cycle
    edges $c_0,c_1,c_2$ are all uncrossed.
\end{lemma}

\begin{proof}
    Again take a disk around the uncrossed curve $ab$ and its endpoints.  The
    four attachment portions have one of the boundary orders

    \[
       0123,\quad 0132,\quad 0231,\quad 0321,
    \]
    according to the two endpoint rotations.  For orders $0123$ and $0321$,
    restore the 4-cycle without a crossing.  For orders $0132$ and $0231$,
    use the single crossing $c_3\mathbin\times s_1$.  The edges $c_0,c_1,c_2$
    remain uncrossed in every case.  The edge $c_3$ is new, and $s_1$ is
    uncrossed outside the disk; hence neither edge has another crossing.

\end{proof}

\begin{lemma}[4-cycle reduction through order 24]\label{lem:order24-reduction}
    Suppose every subcubic graph of order at most 22 is $3$-flexible and
    every biconnected, nonplanar cubic graph of order 24 and girth at least
    five is $2$-flexible.  Then every cubic graph of order at most 24 is
    $2$-flexible.
\end{lemma}

\begin{proof}
    By \cref{lem:flexibility-core}, it is enough to consider a biconnected
    triangle-free cubic graph $G$ of order 24 and a prescribed set $F$ of at
    most two edges.  We may assume that $G$ is nonplanar.  If $G$ has no
    4-cycle, the second hypothesis applies.  Otherwise choose a 4-cycle $C$.

    First suppose that $F$ is not a pair of adjacent edges of $C$.  Choose
    one of the two 4-cycle reductions so that every prescribed cycle
    edge is created inside an endpoint disk: one split handles the opposite
    pair $c_0,c_2$, and the other handles $c_1,c_3$.  In the resulting
    order-22 graph prescribe the new middle edge and the images of all edges
    of $F$ outside $C$.  This is a set of at most three edges.  A drawing
    supplied by $3$-flexibility expands by \cref{lem:4-cycle-expansion};
    prescribed outside edges are unchanged, and prescribed cycle edges are
    drawn uncrossed inside the endpoint disks.

    The remaining case is $F=\{c_i,c_{i+1}\}$.  Rotate the labels so that
    these are $c_0,c_1$, use the displayed split, and prescribe only $ab$ and
    the image of $s_1$.  By \cref{lem:adjacent-4-cycle}, the resulting
    $3$-flexible order-22 drawing expands while keeping both prescribed edges
    uncrossed.
\end{proof}

Consequently, it suffices to verify $3$-flexibility for biconnected,
triangle-free cubic graphs of order at most 22 and $2$-flexibility for
biconnected, nonplanar cubic graphs of order 24 and girth at least five.

\begin{lemma}[5-cycle expansion]\label{lem:5-cycle-expansion}
    Let $C=(v_0,\ldots,v_4)$ be a chordless 5-cycle in a cubic graph $G$,
    and let $x_i$ be the neighbor of $v_i$ outside $C$.  Suppose that the
    vertices $x_0,\ldots,x_4$ are distinct.  Form $H$ by deleting $C$,
    introducing a path $a b c$, and adding the edges
    $ax_0,ax_1,bx_2,cx_3,cx_4$.  If $H$ has a 1-planar drawing in which the
    three edges incident with $b$ are uncrossed, then $G$ is 1-planar.
\end{lemma}

\begin{proof}
    Take a sufficiently small regular neighborhood of the drawn path
    $a b c$.  It contains no crossing and meets the five attachment edges in
    five boundary points.  The rotations at $a,b,c$ give, up to reflection,
    exactly the four boundary orders in \cref{tab:5-cycle-expansion}.  Inside
    the neighborhood, replace the path by the 5-cycle and its five attachment
    portions using the corresponding row.  Here
    $c_i=v_iv_{i+1}$ and $s_i=v_ix_i$, with indices modulo five.

    Every row is a matching of local edges.  The only old edge that occurs in
    a listed crossing is $s_2$, which replaces the uncrossed edge $bx_2$.
    Every new cycle edge is crossed at most once, and all crossings outside
    the neighborhood are unchanged.  Planarizing each displayed crossing
    gives a planar graph in a disk with the stated boundary order and
    alternating incidences at its crossing vertices.  Thus the four rows
    give valid local 1-planar drawings and complete the expansion.
\end{proof}

\begin{table}[!ht]
    \centering
    \caption{Crossing pairs for the 5-cycle expansion.}
    \label{tab:5-cycle-expansion}
    \begin{tabular}{@{}>{$}c<{$}@{\qquad}>{$}l<{$}@{}}
        \toprule
        \multicolumn{1}{c}{boundary order} & \multicolumn{1}{c}{crossing pairs} \\
        \midrule
        01234&\varnothing\\
        01243&c_2\mathbin\times c_4\\
        01342&c_4\mathbin\times s_2\\
        01432&c_1\mathbin\times c_4\\
        \bottomrule
    \end{tabular}
\end{table}

\begin{lemma}[Overlapping 5-cycle expansion]
    \label{lem:overlapping-5-cycle-expansion}
    Use the notation of \cref{lem:5-cycle-expansion}, with $v_2$ assigned
    to the middle vertex $b$.  Suppose that the path $abc$ lies on another
    5-cycle $D$ of $H$.  Replace $D$ by a three-vertex path, assigning the
    attachment at $b$ to its middle vertex, and call the resulting order-24
    graph $K$.  Suppose that this reduction is simple and connected.  If
    $K$ has a 1-planar drawing in which the two edges of the new path are
    uncrossed, then $G$ is 1-planar.
\end{lemma}

\begin{proof}[Proof (computational)]
    The two vertices of $D-\{a,b,c\}$ are outside the original 5-cycle.
    Girth five excludes the case in which they attach at consecutive
    vertices of the original cycle.  Up to a dihedral relabeling, the patch
    to be restored therefore consists of
    $v_0,\ldots,v_4,x_0,x_3$, with the additional edges
    $v_0x_0,v_3x_3,x_0x_3$.

    Take a regular neighborhood of the uncrossed path in $K$.  Its five
    attachment portions occur in the same four boundary orders as in
    \cref{tab:5-cycle-expansion}.  The local drawings in
    \cref{tab:overlapping-5-cycle-expansion} restore the seven-vertex patch.
    Here $d_0=v_0x_0$ and $d_3=v_3x_3$.  Every displayed crossing uses two
    edges of the restored patch, so no further edge of the drawing of $K$
    must be kept uncrossed.  As for the other local tables, the computation
    planarized every displayed crossing and checked planarity in a disk with
    the prescribed boundary order.
\end{proof}

\begin{table}[!ht]
    \centering
    \caption{Crossing pairs for the overlapping 5-cycle expansion.}
    \label{tab:overlapping-5-cycle-expansion}
    \begin{tabular}{@{}>{$}c<{$}@{\qquad}>{$}l<{$}@{}}
        \toprule
        \multicolumn{1}{c}{boundary order} &
        \multicolumn{1}{c}{crossing pairs} \\
        \midrule
        01234&c_4\mathbin\times d_3\\
        01243&\{c_2\mathbin\times c_4,d_0\mathbin\times d_3\}\\
        01342&c_3\mathbin\times d_0\\
        01432&c_2\mathbin\times d_0\\
        \bottomrule
    \end{tabular}
\end{table}

\begin{lemma}[Contracted 5-cycle expansion]\label{lem:5-cycle-star-expansion}
    Let $C=(v_0,\ldots,v_4)$ be a chordless 5-cycle in a cubic graph $G$,
    and let $x_i$ be the neighbor of $v_i$ outside $C$.  Suppose that the
    vertices $x_0,\ldots,x_4$ are distinct, and let $J$ be obtained by
    contracting $C$ to one vertex.  Write $c_i=v_iv_{i+1}$ and
    $s_i=v_ix_i$, with indices modulo five, and call the edges $s_i$ the
    attachments of $C$.
    \begin{enumerate}[(i)]
        \item If $J$ has a 1-planar drawing in which all five edges
        $s_0,\ldots,s_4$ are uncrossed, then, for every $i$, the graph $G$
        has a 1-planar drawing in which the three edges incident with $v_i$
        are uncrossed.
        \item Let $e\in E(J)$.  If $e$ is not an attachment, suppose that
        $J$ has a 1-planar drawing in which $e$ and at least three
        attachments are uncrossed.  If $e=s_i$, instead suppose that $e$ and
        at least three of the other four attachments are uncrossed.  In
        either case, $G$ has a 1-planar drawing in which $e$ is uncrossed.
    \end{enumerate}
\end{lemma}

\begin{proof}[Proof (computational)]
    Take a disk around the contracted vertex.  Fixing $s_0$ as the first
    attachment and identifying an order with its reversal leaves twelve
    possible cyclic orders on the boundary of the disk.  The second column
    of \cref{tab:5-cycle-star-expansion} gives a local drawing for each order
    whose attachment crossings use only $s_0,s_1$.  Relabeling the cycle maps
    any chosen consecutive pair $s_j,s_{j+1}$ to this pair.  Any three
    attachments contain a consecutive pair.  If $e$ is an attachment, any
    three of the other four contain a consecutive pair avoiding $e$.
    Choosing this pair ensures that the local drawing does not cross $e$;
    if $e$ is not an attachment, it lies outside the replacement disk.
    This proves (ii).

    For (i), it is enough by symmetry to preserve the star
    $\{c_4,c_0,s_0\}$.  The third column gives a local drawing for every
    boundary order in which these three edges are uncrossed.  Its attachment
    crossings use only $s_1,s_2,s_3,s_4$, all of which are uncrossed outside
    the disk.  Each listed set of crossings is a matching.  Replacing every
    listed crossing by a degree-four vertex gives a planar graph in the disk
    with the prescribed boundary order and alternating continuations at each
    crossing.  Thus every listed local drawing is 1-planar, and inserting it
    into the disk proves both statements.
\end{proof}

\begin{table}[!ht]
    \centering
    \caption{Crossing pairs for the 5-cycle star expansion.  The ordinary
    column permits local crossings on $s_0,s_1$; the last column preserves
    the star $\{c_4,c_0,s_0\}$.}
    \label{tab:5-cycle-star-expansion}
    \begin{tabular}{@{}>{$}c<{$}@{\qquad}>{$}l<{$}@{\qquad}>{$}l<{$}@{}}
        \toprule
        \multicolumn{1}{c}{boundary order} &
        \multicolumn{1}{c}{ordinary expansion} &
        \multicolumn{1}{c}{star at $v_0$} \\
        \midrule
        01234&\varnothing&\varnothing\\
        01243&c_2\mathbin\times c_4&c_2\mathbin\times s_4\\
        01324&c_1\mathbin\times c_3&c_1\mathbin\times c_3\\
        01342&\{c_1\mathbin\times s_0,c_2\mathbin\times c_4\}&
               \{c_1\mathbin\times c_3,s_2\mathbin\times s_4\}\\
        01423&\{c_1\mathbin\times c_3,c_4\mathbin\times s_1\}&
               c_1\mathbin\times s_4\\
        01432&c_1\mathbin\times c_4&
               \{c_1\mathbin\times s_4,c_3\mathbin\times s_2\}\\
        02134&c_0\mathbin\times c_2&c_2\mathbin\times s_1\\
        02143&c_2\mathbin\times s_0&
               \{c_2\mathbin\times s_1,s_3\mathbin\times s_4\}\\
        02314&c_3\mathbin\times s_1&c_3\mathbin\times s_1\\
        02413&\{c_1\mathbin\times c_4,c_2\mathbin\times s_0,
                 c_3\mathbin\times s_1\}&
               \{c_1\mathbin\times c_3,s_1\mathbin\times s_3,
                 s_2\mathbin\times s_4\}\\
        03124&\{c_0\mathbin\times c_2,c_3\mathbin\times s_0\}&
               \{c_1\mathbin\times c_3,s_1\mathbin\times s_3\}\\
        03214&c_0\mathbin\times c_3&
               \{c_1\mathbin\times s_3,c_3\mathbin\times s_1\}\\
        \bottomrule
    \end{tabular}
\end{table}

\section{Upper Bound}

Our goal is to prove that every simple cubic graph with at most 28 vertices
is 1-planar.  Testing all graphs directly is impractical: there are
$2{,}094{,}480{,}864$ connected cubic graphs with $n=26$ and
$40{,}497{,}138{,}011$ with $n=28$.  Even at ten milliseconds per graph, the
latter family alone would require more than twelve CPU-years.

The reductions above leave five computational claims.  All enumerations below
are exhaustive and contain one representative of each isomorphism class.

We first describe the mathematical objects checked in these computations.  A
\emph{crossing pattern} in a graph $G$ is a set $M$ of unordered pairs of
nonadjacent edges such that no edge occurs in more than one pair.  Let
$\operatorname{supp}(M)$ be the set of edges occurring in the pairs of $M$,
and let $G_M$ be obtained by replacing each pair $\{uv,xy\}\in M$ with a new
vertex adjacent to $u,v,x,y$.

If $G_M$ is planar, then $G$ is 1-planar.  Indeed, consider a planar embedding
of $G_M$ and a small disk around each new vertex.  If the two portions of each
original edge alternate around the new vertex, replace that vertex by a
crossing; otherwise join the two portions of each original edge without a
crossing.  Since no edge belongs to two pairs of $M$, every edge is crossed at
most once.  Moreover, every edge outside $\operatorname{supp}(M)$ is
uncrossed.  The computational proofs below consist of finding such crossing
patterns and testing the planarity of the corresponding graphs $G_M$.  For a
planar graph, the empty crossing pattern suffices.

\begin{claim}\label{clm:flexible22}
    Every connected cubic graph with $n \leq 18$ is $3$-flexible.  Every
    biconnected, nonplanar, triangle-free cubic graph with
    $n \in \{20,22\}$ is $3$-flexible.
\end{claim}

\begin{proof}[Proof (computational)]
    These families contain $1{,}575{,}835$ isomorphism classes.  For every
    nonplanar graph $G$ in the enumeration and every set
    $F\subseteq E(G)$ with $|F|\leq 3$, the computation found a crossing
    pattern $M$ such that $G_M$ is planar and
    $F\cap\operatorname{supp}(M)=\varnothing$.  Hence $G$ has a 1-planar
    drawing in which every edge of $F$ is uncrossed.
\end{proof}

\begin{claim}\label{clm:flexible24}
    Every biconnected, nonplanar cubic graph with $n=24$ and girth $\geq 5$
    is $2$-flexible.
\end{claim}

\begin{proof}[Proof (computational)]
    This family contains $1{,}620{,}470$ isomorphism classes.  For every graph
    $G$ in the enumeration and every set $F\subseteq E(G)$ with $|F|\leq 2$,
    the computation found a crossing pattern $M$ such that $G_M$ is planar and
    $F\cap\operatorname{supp}(M)=\varnothing$.  Thus every edge of $F$ is
    uncrossed in the resulting 1-planar drawing.
\end{proof}

\begin{claim}\label{clm:flexible26}
    Every biconnected, nonplanar cubic graph with $n=26$ and girth
    at least five is $1$-flexible.
\end{claim}

\begin{proof}[Proof (computational)]
    First consider the graphs of girth five.  For each such graph $G$, the
    computation contracted every 5-cycle, canonically labeled the resulting
    order-22 graph, and retained the lexicographically least code.  Equal
    retained graphs were verified only once.  Write $J$ for one retained
    graph and $s_0,\ldots,s_4$ for the five edges incident with its
    degree-five vertex.

    The computation found a crossing pattern with all five $s_i$ uncrossed.
    In addition, for every edge $e\in E(J)$ it found a crossing pattern in
    which $e$ and at least three suitable attachments are uncrossed: all five
    attachments are eligible when $e$ is not one of them, and the other four
    are eligible when $e$ is an attachment.  Each resulting planarization was
    checked.  These conditions do not depend on the cyclic order at the
    contracted vertex.  An edge on the contracted 5-cycle is preserved using
    part (i) of \cref{lem:5-cycle-star-expansion}; every other edge is
    preserved using part (ii).  Thus every parent represented by $J$ is
    $1$-flexible.

    The remaining order-26 graphs have girth at least six.  For every such
    graph $G$ and every edge $e\in E(G)$, the computation directly found a
    crossing pattern whose planarization is planar and whose support avoids
    $e$.
\end{proof}

\begin{claim}\label{clm:girth5-28}
    Every biconnected, nonplanar cubic graph with $n=28$ and girth $=5$ is
    1-planar.
\end{claim}

\begin{proof}[Proposed computational proof]
    For every graph $G$ in this family, the preparation enumerates the
    5-cycle-to-path reductions $(H,b)$ from
    \cref{lem:5-cycle-expansion}.  Girth five makes every first reduction
    simple.  Write the replacement path as $a b c$.

    If the replacement path lies on a second 5-cycle, the preparation
    performs the second path reduction with $b$ assigned to its middle
    attachment.  When the resulting order-24 graph $K$ is biconnected and
    has girth at least five, it is either planar or covered by
    \cref{clm:flexible24}.  Thus its two path edges can be kept uncrossed, and
    \cref{lem:overlapping-5-cycle-expansion} proves that $G$ is 1-planar.

    Otherwise the preparation finds one of two overlapping cycles.  With
    $L,M$ denoting attachments at $a,b$, respectively, their cyclic words are
    \[
        b\,a\,L\,O\,M
        \qquad\hbox{or}\qquad
        b\,a\,L\,O\,O\,M .
    \]
    Contract the complete union of the original 5-cycle and this overlap to
    one vertex.  The resulting quotient has order 21 and one vertex of
    degree six in the first case, or order 20 and one vertex of degree seven
    in the second case; every other vertex has degree three.  Equal
    canonical quotients are checked only once.

    The finite local check in \texttt{cubic/test\_hub\_expansions.py} gives
    the following exact expansion rule.  Every one of the 60 boundary orders
    of the degree-six patch expands.  For the degree-seven patch, 355 of the
    360 boundary orders expand; the five exceptions are
    \[
      0245136,\quad0246135,\quad0264153,\quad
      0315426,\quad0316425.
    \]
    The check surrounds the patch by an uncrossed wheel in the prescribed
    order and verifies the resulting crossing pattern and planarization.

    For a degree-six quotient, the computation finds a drawing in which all
    six hub edges are uncrossed.  For a degree-seven quotient, it considers
    every bijection between the seven patch attachments and the seven hub
    edges.  For each bijection it finds a drawing with all hub edges
    uncrossed whose rotation avoids the five exceptional orders.  Drawings
    are shared among all bijections they cover.  The rotation is read from
    the verified two-page embedding of the planarization, so it is part of
    the certificate rather than a postulated embedding choice.  Replacing
    the hub by the corresponding local drawing restores $G$.

    The preparation fails unless every input graph obtains one of these
    three certificates.  Thus a completed run proves the claim without using
    Claim~3.  The global preparation and verification described in
    \texttt{cubic/verification.md} have not yet been run, so this proposed
    computation does not yet prove the claim.
\end{proof}

\begin{claim}\label{clm:girth6-28}
    Every biconnected, nonplanar cubic graph with $n=28$ and girth $\geq 6$ is
    1-planar.
\end{claim}

\begin{proof}[Proof (computational)]
    This family contains $4{,}624{,}501$ isomorphism classes.  For every graph
    $G$ in the enumeration, the computation found a crossing pattern $M$ such
    that $G_M$ is planar.  Hence every graph in the family is 1-planar.
\end{proof}

The running times for the five computations are summarized in
\cref{tab:cubic28-runtime}.  Claims 1, 2, and 5 are completed computations.
The earlier marked-star preparation for \cref{clm:girth5-28} produced
$38{,}120{,}815$ groups whose measured verification would take 392--470
hours; that route was abandoned.

For Claim~3, a bounded parent-weighted pilot sampled 300 parents from 256
spread-out fine Minibaum parts containing 1,370,731 parents.  Enumerating the
complete reverse fiber of each selected contraction gave
\(\mathop{\rm E}(1/f)=0.1450\), projecting 4.54 million distinct cores.
Timing the same cores gave
\(\mathop{\rm E}(t/f)=0.06790\) seconds per parent; its 95\% sampling interval
was \(0.05887\)--\(0.07693\).  Thus the girth-five part of Claim~3 pilot
projected 19.7--23.6 hours at the measured 30--25-fold effective speedup,
with a 17.1--26.8 hour sampling envelope.  This pilot did not bound the
tail: the production run subsequently encountered long-running cores that
were absent from the sample.  Claim~3 is therefore now run with a
five-second solver limit per core.  Timed-out cores are retained in a
residue for separate verification, and the pilot interval is no longer used
as a completion-time forecast.

For Claim~4, a conservative 60,000-parent structural sample contained 20,056 direct
Claim~2 certificates, 28,391 degree-six overlaps, and 11,553 degree-seven
overlaps.  The final orientation-complete selector can reclassify some
degree-seven cases as the faster degree-six type, so these weights do not
underestimate the verification cost.  A bounded parent-weighted pilot then selected ten quotients of
each overlap type and enumerated their complete reverse fibers.  The
degree-six fibers ranged from 47 to 328 and the degree-seven fibers from 18
to 458.  Timing the same quotients gave
\(\mathop{\rm E}(t/f)=0.000395\) and \(0.001196\) seconds per parent,
respectively.  Weighting these values by the structural sample projects
about 76 CPU-hours for Claim~4 verification.  The complete preparation is
expected to remain the larger Claim~4 term; the previous server preparation
measurements took 15--24 wall-clock hours.

For the production budget we use the largest observed matched value of
\(t/f\) in each overlap class for every parent, rather than the sample mean.
This gives 298 CPU-hours, or about 12 hours at the measured 25-fold effective
speedup.  Together with the 15--24-hour Claim~4 preparation range, this
still projects 27--36 hours for Claim~4.  There is no reliable joint
Claims~3--4 estimate until the bounded Claim~3 run has measured its hard
residue.

\begin{table}[!ht]
    \small
    \centering
    \caption{Measured and projected wall-clock times on 48 cores.}
    \label{tab:cubic28-runtime}
    \begin{tabular}{@{}lrrc@{}}
        \toprule
        computation & instances & running time & verified \\
        \midrule
        \cref{clm:flexible22} & $1{,}575{,}835$ & 14 hours 41 minutes (measured) & yes \\
        \cref{clm:flexible24} & $1{,}620{,}470$ & 41 hours 21 minutes (measured) & yes \\
        \cref{clm:flexible26} & $4.54$ million${}^\ast$ & bounded pass in progress & no \\
        \cref{clm:girth5-28}  & $5.3$ million${}^\dagger$ & 27--36 hours (projected) & no \\
        \cref{clm:girth6-28}  & $4{,}624{,}501$ & 3 days (measured) & yes \\
        \midrule
        remaining Claims 3--4 & & not yet estimated & no \\
        \bottomrule
    \end{tabular}
    \par\smallskip
    {\footnotesize $^\ast$Projected number of distinct Claim 3 cores.
    $^\dagger$Projected number of distinct degree-six and degree-seven
    quotients; the running time includes Claim~4 preparation.\par}
\end{table}

Thus the present computations do not establish
\cref{clm:flexible26,clm:girth5-28}, and the theorem below remains
conditional on completing or
replacing those two computations.

\begin{theorem}\label{thm:cubic28}
    Every simple cubic graph with at most 28 vertices is 1-planar.
\end{theorem}

\begin{proof}
    A planar graph is $k$-flexible for every $k$, so every obstruction to a
    flexibility statement below is nonplanar.  We first apply the flexibility
    results at the smaller orders.  If a subcubic graph of order at most 22
    were not $3$-flexible, a smallest such graph would be biconnected, cubic,
    and triangle-free by
    \cref{lem:flexibility-core}, contrary to \cref{clm:flexible22}.  Thus every
    subcubic graph of order at most 22 is $3$-flexible.  It follows from
    \cref{lem:order24-reduction,clm:flexible24} that every cubic graph of order
    at most 24 is $2$-flexible.  A smallest subcubic graph of order at most 24
    that is not $2$-flexible would again be cubic by
    \cref{lem:flexibility-core}; consequently, every subcubic graph of order
    at most 24 is $2$-flexible.

    We next show that every subcubic graph of order at most 26 is
    $1$-flexible.  Suppose otherwise, and choose a smallest counterexample
    $G$.  By \cref{lem:flexibility-core}, the graph $G$ is biconnected, cubic,
    and triangle-free.  The preceding paragraph excludes orders at most 24,
    and a cubic graph has even order; hence $G$ has order 26.  Fix an edge
    $e\in E(G)$.  If $G$ contains a 4-cycle, choose one of the
    two reductions in \cref{lem:4-cycle-expansion} so that, when $e$ belongs
    to the cycle, it is restored inside an endpoint disk.  In the resulting
    graph of order 24, prescribe the edge $ab$ and, when $e$ lies outside the
    cycle, its image.  Its $2$-flexibility gives a drawing that expands to a
    drawing of $G$ in which $e$ is uncrossed.  If $G$ has no 4-cycle, then it
    has girth at least five, so \cref{clm:flexible26} directly gives a
    1-planar drawing in which $e$ is uncrossed.  Since $e$ was arbitrary,
    both cases contradict the choice of $G$.

    Finally, suppose that a subcubic graph of order at most 28 is not
    1-planar, and choose a smallest such graph $G$.  By
    \cref{lem:flexibility-core} with $k=0$, the graph $G$ is biconnected,
    cubic, and triangle-free.  The preceding paragraph excludes orders at
    most 26, and cubic graphs have even order; hence $G$ has order 28.  If $G$
    contains a 4-cycle, apply
    \cref{lem:4-cycle-expansion}.  The resulting graph has order 26, and its
    $1$-flexibility provides a 1-planar drawing in which $ab$ is uncrossed;
    this drawing expands to a 1-planar drawing of $G$.  Otherwise $G$ has
    girth at least five.  If its girth is five, \cref{clm:girth5-28} applies;
    if its girth is at least six, \cref{clm:girth6-28} applies.  Every case is
    a contradiction.
\end{proof}





\bibliographystyle{plainurl}
\bibliography{cubic}

\end{document}
